\documentclass[11pt,a4paper]{article}
\usepackage[margin=1in]{geometry}
\usepackage[T1]{fontenc}
\usepackage[utf8]{inputenc}
\usepackage{lmodern}
\usepackage{setspace}
\usepackage{booktabs}
\usepackage{hyperref}
\usepackage{enumitem}
\usepackage{titlesec}

\setstretch{1.05}
\setlength{\parskip}{0.5em}
\setlength{\parindent}{0pt}
\titleformat{\section}{\large\bfseries}{\thesection.}{0.5em}{}
\titleformat{\subsection}{\normalsize\bfseries}{\thesubsection}{0.5em}{}

\title{Confidence-Aware ProSST: Improving Robustness of Structure-Aware Protein Language Models Under Structural Uncertainty}
\author{
    Member 1 (NTU Email) \and
    Member 2 (NTU Email) \and
    Member 3 (NTU Email)
}
\date{AI6127 Final Project Report}

\begin{document}
\maketitle

\begin{abstract}
This project studies robust protein language modeling under structural uncertainty. We extend ProSST with confidence-aware sequence-structure integration using residue-level confidence signals (for example pLDDT). We compare hard masking, soft weighting, and learned gating strategies in disentangled attention. On ProteinGym, our best variant aims to improve robustness in low-confidence structure regimes while maintaining competitive overall performance against ProSST baselines. These results suggest that uncertainty-aware structure fusion is important for practical deployment of structure-aware protein language models.
\end{abstract}

\section{Introduction}
\subsection{Problem and Motivation}
Protein language models can benefit from structure signals, but predicted structures have non-uniform quality across proteins and residues. This report explores how to integrate confidence signals so the model trusts structure when reliable and relies more on sequence when uncertain.

\subsection{Challenges}
\begin{itemize}[leftmargin=*]
    \item \textbf{Data quality}: Structural confidence varies by residue.
    \item \textbf{Modeling}: Sequence-structure fusion should be confidence-aware.
\end{itemize}

\subsection{Contributions}
\begin{itemize}[leftmargin=*]
    \item A confidence-aware extension of ProSST.
    \item Three variants: hard masking, soft weighting, learned gating.
    \item Stratified robustness analysis by confidence bins.
\end{itemize}

\section{Related Work / Background}
\subsection{Protein Language Models}
Briefly review sequence-only and structure-aware PLMs (for example ESM, SaProt, ProSST), and discuss their assumptions.

\subsection{Confidence and Uncertainty}
Summarize reliability-aware modeling ideas and motivate confidence-guided attention.

\subsection{Positioning}
Explain how this work differs from ProSST and why confidence-aware fusion is a practical extension.

\section{Approach}
\subsection{Baseline: ProSST}
Describe structure tokenization and sequence-structure disentangled attention.

\subsection{Proposed Method: CA-ProSST}
\textbf{V1 Hard Masking:} mask low-confidence structure tokens. \\
\textbf{V2 Soft Weighting:} scale structure-related attention by confidence. \\
\textbf{V3 Learned Gate:} learn a gate per residue to balance sequence and structure contributions.

\subsection{Implementation Details}
State architecture changes, confidence feature design, and training details.

\section{Experiments}
\subsection{Data}
Describe datasets, references, and exact task input/output format.

\subsection{Evaluation Metrics}
Define primary and secondary metrics (for example Spearman, NDCG, Top-recall).

\subsection{Experimental Settings}
Report hardware, optimizer, learning rate, batch size, training steps, and random seed policy.

\subsection{Baselines}
Include ProSST, sequence-only ablations, and optional external baselines.

\subsection{Main Results}
\begin{table}[h]
\centering
\caption{Overall performance placeholder}
\begin{tabular}{lccc}
\toprule
Model & Spearman & NDCG & Top-recall \\
\midrule
ProSST baseline & -- & -- & -- \\
CA-ProSST V1 & -- & -- & -- \\
CA-ProSST V2 & -- & -- & -- \\
CA-ProSST V3 & -- & -- & -- \\
\bottomrule
\end{tabular}
\end{table}

\subsection{Ablations and Robustness}
Report component ablations and confidence-stratified performance.

\subsection{Result Commentary}
Comment on whether outcomes match expectations, and explain why.

\section{Analysis}
Include qualitative analysis: success cases, failure cases, and model behavior.

\section{Conclusion}
Summarize findings, lessons learned, limitations, and future work.

\section*{Team Contributions}
\begin{itemize}[leftmargin=*]
    \item Member 1 (xx\%): 1--2 sentence contribution summary.
    \item Member 2 (xx\%): 1--2 sentence contribution summary.
    \item Member 3 (xx\%): 1--2 sentence contribution summary.
\end{itemize}

\section*{References}
Add all cited references here.

\appendix
\section{Optional Appendix}
Extra tables, figures, examples, and implementation details can be added here.

\end{document}
